\subsection*{1. Deriving the Expected Returns and Covariance Matrix}
We are given the riskless rate $R^0 = 0.06$ and the relationship between the two risky assets:
\[
\rho(S^2) = 0.06 + 0.2\rho(S^1)
\]
The probability distribution for $\rho(S^1)$ is provided as follows:
\begin{table}[H]
\centering
\begin{tabular}{lccccc}
\toprule
$\rho(S^1)$ \& Probability \\
\midrule
0.6 \& 0.1 \\
0.4 \& 0.2 \\
0.3 \& 0.4 \\
0.2 \& 0.2 \\
-1  \& 0.1 \\
\bottomrule
\end{tabular}
\end{table}

\vspace{0.5em}
\noindent \textbf{Expected Returns:}
First, we calculate the expected return for the first asset, $\mathbb{E}[\rho(S^1)]$:
\[
\mathbb{E}[\rho(S^1)] = (0.6)(0.1) + (0.4)(0.2) + (0.3)(0.4) + (0.2)(0.2) + (-1)(0.1)
\]
\[
\mathbb{E}[\rho(S^1)] = 0.06 + 0.08 + 0.12 + 0.04 - 0.10 = 0.2
\]
Using the linear relationship, we can directly find the expected return for the second asset, $\mathbb{E}[\rho(S^2)]$:
\[
\mathbb{E}[\rho(S^2)] = 0.06 + 0.2\mathbb{E}[\rho(S^1)] = 0.06 + 0.2(0.2) = 0.06 + 0.04 = 0.1
\]
Thus, the expected return vector is $R = \begin{pmatrix} 0.2 \\ 0.1 \end{pmatrix}$.

\vspace{0.5em}
\noindent \textbf{Covariance Matrix ($\Sigma$):}
Next, we need the variance of $\rho(S^1)$. We first find $\mathbb{E}[\rho(S^1)^2]$:
\[
\mathbb{E}[\rho(S^1)^2] = (0.6^2)(0.1) + (0.4^2)(0.2) + (0.3^2)(0.4) + (0.2^2)(0.2) + (-1^2)(0.1)
\]
\[
\mathbb{E}[\rho(S^1)^2] = 0.036 + 0.032 + 0.036 + 0.008 + 0.10 = 0.212
\]
Now, calculate the variance:
\[
V(S^1) = \mathbb{E}[\rho(S^1)^2] - (\mathbb{E}[\rho(S^1)])^2 = 0.212 - (0.2)^2 = 0.212 - 0.04 = 0.172
\]
Using the properties of variance and covariance with the linear relationship $\rho(S^2) = 0.06 + 0.2\rho(S^1)$:
\[
V(S^2) = V(0.06 + 0.2\rho(S^1)) = 0.2^2 V(S^1) = 0.04 \times 0.172 = 0.00688
\]
\[
Cov(S^1, S^2) = Cov(\rho(S^1), 0.06 + 0.2\rho(S^1)) = 0.2 V(S^1) = 0.2 \times 0.172 = 0.0344
\]
This gives us the exact covariance matrix provided in the exercise:
\[
\Sigma = \begin{pmatrix} 0.172 \& 0.0344 \\ 0.0344 \& 0.00688 \end{pmatrix}
\]
Notice that the determinant of $\Sigma$ is zero ($0.172 \times 0.00688 - 0.0344^2 = 0$). The assets are perfectly positively correlated, meaning $\Sigma$ is not invertible and we cannot use the standard matrix optimization from Exercise 1.

\subsection*{2. Minimum-Variance Portfolio (Long \& Short allowed)}
We want to minimize portfolio variance $V(w)$ subject to $w_1 + w_2 = 1$. Let's express the portfolio return $V^w$ strictly in terms of $w_1$ and $\rho(S^1)$.

Substitute $w_2 = 1 - w_1$ and $\rho(S^2) = 0.06 + 0.2\rho(S^1)$ into the portfolio return equation:
\[
V^w = w_1\rho(S^1) + w_2\rho(S^2) = w_1\rho(S^1) + (1 - w_1)(0.06 + 0.2\rho(S^1))
\]
Expand and factor by $\rho(S^1)$:
\[
V^w = w_1\rho(S^1) + 0.06 - 0.06w_1 + 0.2\rho(S^1) - 0.2w_1\rho(S^1)
\]
\[
V^w = (w_1 + 0.2 - 0.2w_1)\rho(S^1) + 0.06(1 - w_1)
\]
\[
V^w = (0.2 + 0.8w_1)\rho(S^1) + 0.06(1 - w_1)
\]
The variance of this portfolio is simply the variance of the random component:
\[
V(w) = Var(V^w) = (0.2 + 0.8w_1)^2 V(S^1) = (0.2 + 0.8w_1)^2 (0.172)
\]
To find the minimum variance, we can just set this squared term to zero (since variance cannot be negative):
\[
0.2 + 0.8w_1 = 0 \implies 0.8w_1 = -0.2 \implies \tilde{w}_1 = -0.25
\]
Since $w_1 + w_2 = 1$, we have:
\[
\tilde{w}_2 = 1 - (-0.25) = 1.25
\]
So the optimal weights are $\tilde{w} = (-0.25, 1.25)^T$. This means we short asset 1 and go long on asset 2.

\vspace{0.5em}
\noindent \textbf{Expected Value \& Variance:}
By construction, the variance of this portfolio is zero:
\[
Var(V) = 0
\]
The expected return is:
\[
r(V) = \tilde{w}^T R = -0.25(0.2) + 1.25(0.1) = -0.05 + 0.125 = 0.075
\]
Notice that the risk-free rate is $R^0 = 0.06$, but we have created a portfolio with absolutely no risk ($Var(V) = 0$) that yields $0.075 > 0.06$. 

\subsection*{3. Creating an Arbitrage Opportunity}
Since we found a risk-free portfolio that pays more than the market risk-free rate, we can construct an arbitrage strategy.

Let $\lambda > 0$ be some amount of capital. We construct our strategy at time $t=0$:
\begin{enumerate}
    \item Borrow $\lambda$ dollars at the risk-free rate $R^0 = 0.06$. (Cash flow: $+\lambda$)
    \item Invest this $\lambda$ into our zero-variance portfolio. Since our weights are $\tilde{w}_1 = -0.25$ and $\tilde{w}_2 = 1.25$:
    \begin{itemize}
        \item Short sell $0.25\lambda$ worth of asset 1. (Cash flow: $+0.25\lambda$)
        \item Buy $1.25\lambda$ worth of asset 2. (Cash flow: $-1.25\lambda$)
    \end{itemize}
\end{enumerate}
Total initial cash flow at $t=0$ is: $\lambda + 0.25\lambda - 1.25\lambda = 0$. We need zero of our own money to start this operation.

At time $t=1$, we liquidate the portfolio and pay back the loan:
\begin{enumerate}
    \item Pay back the borrowed money with interest: $-\lambda(1 + R^0) = -\lambda(1.06)$.
    \item Repurchase asset 1 to close the short position: $-0.25\lambda(1 + \rho(S^1))$.
    \item Sell asset 2 from our long position: $+1.25\lambda(1 + \rho(S^2))$.
\end{enumerate}
Let's calculate the total payoff at $t=1$, substituting $\rho(S^2) = 0.06 + 0.2\rho(S^1)$:
\[
\text{Payoff} = -\lambda(1.06) - 0.25\lambda(1 + \rho(S^1)) + 1.25\lambda(1 + 0.06 + 0.2\rho(S^1))
\]
\[
\text{Payoff} = \lambda \left[ -1.06 - 0.25 - 0.25\rho(S^1) + 1.25(1.06) + 1.25(0.2)\rho(S^1) \right]
\]
\[
\text{Payoff} = \lambda \left[ -1.31 - 0.25\rho(S^1) + 1.325 + 0.25\rho(S^1) \right]
\]
The $\rho(S^1)$ terms perfectly cancel out, mathematically proving our risk is strictly zero.
\[
\text{Payoff} = \lambda(-1.31 + 1.325) = 0.015\lambda
\]
Because $0.015\lambda$ is strictly positive and guaranteed regardless of what the market does, we have mathematically proven an arbitrage opportunity. The larger the $\lambda$ we borrow, the larger our risk-free profit becomes.